%% kapa: exam=1 round=19 year=2008 period=1 form=A subject=appraisal_law subject_ko=부동산관계법규 seq=19 number=59
%% source: Q-Net 감정평가사 기출문제 (cst00309) artlSeq=1001962 file=extracted/1008/1001962/2008년 제19회 감정평가사 1차시험 기출문제/1교시.hwp
%% conversion: kordoc 4.24.51 -> tools/gen1.py (auto); review=proofread-sonnet; flags=none
\begin{question}{59}
부동산 가격공시 및 감정평가에 관한 법령상 벌칙에 관한 설명으로 옳은 것은?
\begin{choices}
  \item 과태료 처분에 불복이 있는 자는 국토해양부장관에게 이의를 제기할 수 있고 이 경우 국토해양부장관은 지체없이 관할법원에 통보하고, 관할법원은 민사소송법에 따라 재판한다.
  \item 감정평가사의 자격증을 다른 사람에게 양도 또는 대여한 자와 이를 양수 또는 대여받은 자는 1년 이하의 징역 또는 1천만원 이하의 벌금에 처한다.
  \item 자기 또는 친족의 소유토지 등 불공정한 감정평가를 할 우려가 있다고 인정되는 토지에 대한 감정평가를 한 경우 1년 이하의 징역 또는 1천만원 이하의 벌금에 처한다.
  \item 과태료를 납부하지 아니한 경우 민사집행법에 따라 징수한다.
  \item 감정평가사의 자격등록이 거부되거나 취소된 자가 토지의 감정평가를 행한 경우 1년 이하의 징역 또는 1천만원 이하의 벌금에 처한다.
\end{choices}
\end{question}
\answer{2}
